\documentclass[]{article}
\usepackage{geometry,tikz,lipsum,amsmath,amssymb,soul,xcolor,cancel}
\usepackage[dvipsnames]{xcolor}
\tikzset{every picture/.style={line width=0.75pt}} %set default line width to 0.75pt        
\numberwithin{equation}{section}

\NewDocumentCommand{\caja}{ m O{gray!40} O{gray!10} m }{
	\begin{figure}[h]
		\centering
		\begin{tikzpicture}
			\node[anchor=text, text width=\columnwidth-1.2cm, draw, rounded corners, line width=1pt, fill=#3, inner sep=5mm] (big) {\\#4};
			\node[draw, rounded corners, line width=.5pt, fill=#2, anchor=west, xshift=5mm] (small) at (big.north west) {#1};
		\end{tikzpicture}
	\end{figure}
}

\newcommand{\deriv}[3]{\left(\frac{\partial #1}{\partial #2} \right)_{#3}}

\newcommand{\dderiv}[3]{\left(\dfrac{\partial #1}{\partial #2} \right)_{#3}}

\NewDocumentCommand{\indUpDown}{ O{\Lambda} O{\mu} O{\nu} }{
	{#1}^{#2}_{\ #3}
}
\NewDocumentCommand{\indDownUp}{ O{\Lambda} O{\mu} O{\nu} }{
	{#1}_{#2}^{\ #3}
}

\newcommand{\LambdaMuNu}{\indUpDown}
\newcommand{\0}{\textcolor{gray}{0}}
\newcommand{\ei}{\mathbf{e}_i}
\newcommand{\ej}{\mathbf{e}_j}
\newcommand{\lc}{\epsilon_{ijk}}
\newcommand{\gr}{\nabla}
\newcommand{\di}{\nabla\cdot}
\newcommand{\ro}{\nabla\times}
\newcommand{\la}{\nabla^2}
\newcommand{\rv}{\mathbf{r}}
\newcommand{\pv}{\mathbf{p}}
\newcommand{\xv}{\mathbf{x}}
\newcommand{\yv}{\mathbf{y}}
\newcommand{\zv}{\mathbf{z}}
\newcommand{\xhat}{\hat{\xv}}
\newcommand{\yhat}{\hat{\yv}}
\newcommand{\zhat}{\hat{\zv}}
\newcommand{\ihat}{\hat{\mathbf{i}}}
\newcommand{\jhat}{\hat{\mathbf{j}}}
\newcommand{\khat}{\hat{\mathbf{k}}}
\newcommand{\rhat}{\hat{\rv}}
\newcommand{\that}{\hat{\mathbf{\theta}}}
\newcommand{\phat}{\hat{\mathbf{\varphi}}}
\newcommand{\rhhat}{\hat{\mathbf{\rho}}}
\newcommand{\nhat}{\hat{\mathbf{n}}}
\newcommand{\Rhat}{\hat{\mathbf{R}}}
\newcommand{\eo}{\epsilon_{0}}
\newcommand{\uo}{\mu_{0}}
\NewDocumentCommand{\ke}{ O{1}}{
	\frac{#1}{4\pi\eo}
}
\NewDocumentCommand{\ku}{ O{}}{
	\frac{\uo #1}{4\pi}
}
%\newcommand{\ke}{\frac{1}{4\pi\eo}}
%\newcommand{\ku}{\frac{\uo}{4\pi}}
\newcommand{\rdif}{\rv - \rv'}
\newcommand{\pd}[2]{\dfrac{\partial #1}{\partial #2}}
\newcommand{\lpd}[2]{\partial #1/\partial #2}
\newcommand{\spd}[2]{\frac{\partial #1}{\partial #2}}
\newcommand{\td}[2]{\dfrac{d #1}{d #2}}
\newcommand{\ltd}[2]{d #1/d #2}
\newcommand{\std}[2]{\frac{d #1}{d #2}}
\newcommand{\n}[1]{\mathbf{#1}}
\newcommand{\eff}[1]{#1_\mathrm{eff}}
\newcommand{\Ueff}{\eff{U}}

\newcommand{\wfrac}[3][3pt]{%
	\frac{\wfracterm{depth}{\dp}{#1}{#2}}{\wfracterm{height}{\ht}{#1}{#3}}%
}
\newcommand{\wfracterm}[4]{%
	\sbox0{$\displaystyle#4$}%
	\vrule width 0pt #1 \dimexpr #20+#3\relax
	\usebox{0}%
}

\newcommand{\pr}[1]{\left\langle #1 \right\rangle }

\newcommand{\qed}{\hfill\(\blacksquare\)}
\newcommand{\resp}{
	\\
	
	\textbf{Respuesta:} }
\newcommand{\celsius}{^\circ\mathrm{C}}

\newcommand{\GeV}{\,\text{GeV}}

%opening
\title{Ejercicios Certamen 2 Mecánica Clásica Avanzada}
\author{Simon Fonseca}

\begin{document}

\maketitle
\section{Problema 1}
A point-like nonrelativistic particle of mass $m$ scatters in the time-dependent spherically symmetric harmonic potential
\begin{equation}
	U(r,t) = \frac{k(t)}{2} r^2, ~~~~~ k(t) = \frac{\kappa}{t^2}, ~~~~~ \kappa > 0
\end{equation}
\begin{enumerate}	
%%%%%%%%%%%%%%%%%%%%%%%%%%%%%%%%%%%%%%%%%%%%%%%%%%%%%%%%%%%%%%%%%%%%%%%%%%%%%%%%%%%%%%%%%%%%%%%%%%%%%%%%
	\item[(a)] Identify all continuous symmetries of the action. Analyze if there are symmetries of type
	\begin{equation}
		\rv \rightarrow \alpha \rv, ~~~~~ t\rightarrow \beta t, ~~~~~ \alpha,\beta=\mathrm{const.}
	\end{equation}
	which do not change the equations of motion.
	\resp El lagrangiano será
	\begin{equation}
		L = \frac{1}{2}m \left(\dot{r}^2 + r^2 \dot{\theta}^2 + r^2\dot{\phi}^2 \sin^2 \theta \right) - \frac{\kappa}{2 t^2} r^2
	\end{equation}
	
	La variable $\phi$ es cíclica, por lo que $M_z$ será conservado. Aun más, el sistema es esféricamente simétrico, por lo que todo el vector de momento angular se conservará:
	\begin{equation}
		\n{M} = \mathrm{const.}
	\end{equation}
	
	Reemplazando $\rv \rightarrow \alpha \rv, t\rightarrow \beta t$ en las derivadas,
	\begin{equation}
		\pd{r}{t} \rightarrow \pd{\alpha r}{\beta t} = \frac{\alpha}{\beta} \pd{r}{t}
	\end{equation}
	Luego en el lagrangiano
	\begin{align*}
		L' &= \frac{1}{2}m \left(\frac{\alpha^2}{\beta^2} \dot{r}^2 + \frac{\alpha^2}{\beta^2} r^2 \dot{\theta}^2 + \frac{\alpha^2}{\beta^2} r^2\dot{\phi}^2 \sin^2 \theta \right) - \frac{\kappa}{2 \beta^2 t^2} \alpha^2 r^2\\
		L' &= \frac{\alpha^2}{\beta^2} \left[\frac{1}{2}m \left(\dot{r}^2 + r^2 \dot{\theta}^2 + r^2\dot{\phi}^2 \sin^2 \theta \right) - \frac{\kappa}{2 t^2} r^2\right]\\
		L' &= \frac{\alpha^2}{\beta^2} L
	\end{align*}
	Ya que multiplicar un lagrangiano por una cantidad constante no cambia las ecuaciones de movimiento, tendremos una simetría sí y solo sí $\alpha,\beta \in \mathbb{R}, \beta\neq 0$.
	
	Sin embargo, al considerar el efecto sobre la acción,
	\begin{equation*}
		S = \int L \, dt ~\rightarrow~ S' = \int L' \, dt' = \int \frac{\alpha^2}{\beta^2} L \, \beta dt = \frac{\alpha^2}{\beta} S
	\end{equation*}
	Para que podamos usar el teorema de Noether la acción deberá ser invariante, por lo que
	\begin{equation}
		\frac{\alpha^2}{\beta} = 1 ~\Rightarrow~ \alpha = \sqrt{\beta}
	\end{equation}
	\begin{equation*}
		\Rightarrow \rv \rightarrow \alpha \rv, ~~~~~ t\rightarrow \alpha^2 t, ~~~~~ \alpha = \mathrm{const.}
	\end{equation*}
	\qed
	
%%%%%%%%%%%%%%%%%%%%%%%%%%%%%%%%%%%%%%%%%%%%%%%%%%%%%%%%%%%%%%%%%%%%%%%%%%%%%%%%%%%%%%%%%%%%%%%%%%%%%%%%
	\item[(b)] Construct the integrals of motion which might exist according to Noether’s theorem. Demonstrate that these quantities are conserved using the equations of motion.
	\resp El momento angular vendrá dado por
	\begin{equation}
		\n{M} = m\rv\times\dot{\rv} = \mathrm{const.}
	\end{equation}
	
	Para la simetría de escala, escribiendo las transformaciones en notación con $\varepsilon$
	\begin{equation*}
		\alpha = e^{\varepsilon(t)} \approx 1 + \varepsilon(t)
	\end{equation*}
	tal que $\alpha$ es generado por cambios infinitesimales $\varepsilon$, los cambios de coordenadas y tiempo será,, entonces
	\begin{equation*}
		\rv' \approx \rv + \delta \rv \approx \rv + \varepsilon \, \rv, ~~~~~ t' \approx t + \delta t \approx t + (1 + \varepsilon)^2 t \approx t + 2\varepsilon \, t
	\end{equation*}
	\begin{equation*}
		\delta \rv \approx \varepsilon \rv, ~~~~~ \delta t \approx 2\varepsilon t
	\end{equation*}
	\begin{align*}
		\Rightarrow d \rv' &= d \left(e^{\varepsilon(t)} \rv\right)\\
		&= e^{\varepsilon} d\rv + \rv e^{\varepsilon} d \varepsilon\\
		&= e^{\varepsilon} \dot{\rv} dt + \rv e^{\varepsilon} \dot{\varepsilon} dt \\
		&\approx (1 + \varepsilon) \dot{\rv} dt + \rv (1 + \varepsilon) \dot{\varepsilon} dt \\
		&\approx \underbrace{\dot{\rv} dt}_{d\rv} + \underbrace{\varepsilon \dot{\rv} dt}_{\mathcal{O}(\varepsilon)} + \underbrace{\rv \dot{\varepsilon} dt}_{\mathcal{O}(\dot{\varepsilon})} + \underbrace{\rv \varepsilon \dot{\varepsilon} dt}_{\mathcal{O}(\varepsilon \dot{\varepsilon})}\\
		&\\
		\Rightarrow d t' &= d \left(e^{2\varepsilon(t)} t\right)\\
		&= e^{2\varepsilon} dt + t e^{2\varepsilon} \cdot 2 d \varepsilon\\
		&= e^{2\varepsilon} dt + 2 t e^{2\varepsilon} \dot{\varepsilon} dt \\
		&\approx (1 + \varepsilon)^2 dt + 2 t (1 + \varepsilon)^2 \dot{\varepsilon} dt \\
		&\approx (1 + 2\varepsilon + \varepsilon^2) dt + 2 t (1 + 2\varepsilon + \varepsilon^2) \dot{\varepsilon} dt \\
		&\approx dt + \underbrace{2\varepsilon dt}_{\mathcal{O}(\varepsilon)} + \underbrace{\varepsilon^2 dt}_{\mathcal{O}(\varepsilon^2)} + \underbrace{2 t \dot{\varepsilon} dt}_{\mathcal{O}(\dot{\varepsilon})} + \underbrace{4 t \varepsilon \dot{\varepsilon} dt}_{\mathcal{O}(\varepsilon \dot{\varepsilon})} + \underbrace{2 t \varepsilon^2 \dot{\varepsilon} dt}_{\mathcal{O}(\varepsilon^2 \dot{\varepsilon})}
	\end{align*}
	Todos los términos $\varepsilon \dot{\varepsilon}, ~ \varepsilon^2, ~ \varepsilon^2 \dot{\varepsilon}$ son insignificantes, y después de esto deberemos diferenciar por $\dot{\varepsilon}$, por lo que solo nos interesarán los términos $\mathcal{O}(\dot{\varepsilon})$
	\begin{equation*}
		d \rv' \approx \left(\dot{\rv} + \dot{\varepsilon} \rv\right) dt, ~~~~~ dt' \approx \left(1 + 2 \dot{\varepsilon} t\right) dt
	\end{equation*}
	Luego, la velocidad
	\begin{equation*}
		\dot{\rv}' \equiv \frac{d \rv'}{d t'} \approx \frac{\left(\dot{\rv} + \dot{\varepsilon} \rv \right) dt}{\left(1 + 2 \dot{\varepsilon} t\right) dt}
	\end{equation*}
	Aproximando la división con \texttt{Mathematica},
	\begin{equation*}
		\dot{\rv}' \approx \dot{\rv} + \delta \dot{\rv} \approx \dot{\rv} + \dot{\varepsilon}\left(\rv - 2 \dot{\rv} t\right)
	\end{equation*}
	\begin{equation*}
		\delta \dot{\rv} \approx \dot{\varepsilon}\left(\rv - 2 \dot{\rv} t\right)
	\end{equation*}
	Luego, aproximando el lagrangiano a primer orden,
	\begin{align*}
		L'(\rv',\dot{\rv}',t') &= L'(\rv + \delta \rv, \dot{\rv} + \delta \dot{\rv}, t + \delta t)\\
		&\approx L(\rv,\dot{\rv},t) + \pd{L}{\rv} \cdot \delta\rv + \pd{L}{\dot{\rv}} \cdot \delta\dot{\rv} + \pd{L}{t} \delta t\\
		&\approx L(\rv,\dot{\rv},t) + \pd{L}{\rv} \cdot \varepsilon \rv + \pd{L}{\dot{\rv}} \cdot \dot{\varepsilon}\left(\rv - 2 \dot{\rv} t\right) + \pd{L}{t} 2 \varepsilon t
	\end{align*}
	De nuevo, solo nos interesarán los términos $\mathcal{O}(\dot{\varepsilon})$, por lo que, considerando que $\lpd{L}{\dot{\rv}} = m \dot{\rv}$
	\begin{equation*}
		L' \approx L + m \dot{\rv} \cdot \dot{\varepsilon}\left(\rv - 2 \dot{\rv} t\right) 
	\end{equation*}
	Por lo que la acción
	\begin{align*}
		S' &= \int L' \, dt'\\
		&\approx \int \left[L + m \dot{\rv} \cdot \dot{\varepsilon}\left(\rv - 2 \dot{\rv} t\right)\right] \, \left[1 + 2 \dot{\varepsilon} t\right] dt\\
		&\approx \int L \, dt + \int m \dot{\rv} \cdot \dot{\varepsilon} \left(\rv - 2 \dot{\rv} t\right) \, dt + \int 2 \dot{\varepsilon} t L \, dt + \int 2 m \dot{\rv} \cdot \dot{\varepsilon}^2 t \left(\rv - 2 \dot{\rv} t\right) \, dt
	\end{align*}
	El último término es $\mathcal{O}(\dot{\varepsilon}^2)$, por lo que es insignificante. Tendremos entonces
	\begin{equation*}
		S' \approx \int \left[L + m \dot{\rv} \cdot \dot{\varepsilon} \left(\rv - 2 \dot{\rv} t\right) + 2 \dot{\varepsilon} t L\right] \, dt
	\end{equation*}
	Y la variación de $L$ será
	\begin{equation*}
		S' \approx \int \left[L + \delta L\right] dt \approx \int \left[L + \dot{\varepsilon} \left( m \dot{\rv} \cdot \left(\rv - 2 \dot{\rv} t\right) + 2 t L\right)\right] \, dt
	\end{equation*}
	\begin{equation*}
		\delta L \approx \dot{\varepsilon}\left[ m \dot{\rv} \cdot \left(\rv - 2 \dot{\rv} t\right) + 2 t L\right]
	\end{equation*}
	Finalmente, la cantidad conservada será
	\begin{align*}
		Q &= \pd{\delta L}{\dot{\varepsilon}}\\
		&= m \dot{\rv} \cdot \left(\rv - 2 \dot{\rv} t\right) + 2 t L\\
		&= \pv \cdot \rv - 2 m \dot{\rv}^2 t + 2 t \left(\frac{1}{2}m \dot{\rv}^2 - U\right)\\
		&= \pv \cdot \rv - 2 t \left(\frac{1}{2}m \dot{\rv}^2 + U\right)\\
		&= \pv \cdot \rv - 2 H t\\
	\end{align*}
	\qed
	
	
%%%%%%%%%%%%%%%%%%%%%%%%%%%%%%%%%%%%%%%%%%%%%%%%%%%%%%%%%%%%%%%%%%%%%%%%%%%%%%%%%%%%%%%%%%%%%%%%%%%%%%%%
	\item[(c)] We have seen in lectures that in static (time-independent) Coulomb potential the particle's trajectory always remains in the plane (so-called scattering plane). Analyze if this property remains valid in the time-dependent potential given in Eq. (1). Analyze if the Kepler’s laws remain valid in this potential	field.
	\resp El potencial es central y el momento angular $\n{M}$ es conservado componente por componente, por lo que el plano de rotación será constante.
	
	La segunda ley de Kepler se conservará porque el campo es central; la primera y tercera no lo harán, ya que dependen de una trayectoria específica que en este caso será alterada.
	
	\qed
\end{enumerate}

%%%%%%%%%%%%%%%%%%%%%%%%%%%%%%%%%%%%%%%%%%%%%%%%%%%%%%%%%%%%%%%%%%%%%%%%%%%%%%%%%%%%%%%%%%%%%%%%%%%%%%%%

\section{Problema 2}
A relativistic particle of mass $m$ moves in the spherically symmetric potential $U (r) = \alpha/r + \beta/r^2$.
\begin{enumerate}
%%%%%%%%%%%%%%%%%%%%%%%%%%%%%%%%%%%%%%%%%%%%%%%%%%%%%%%%%%%%%%%%%%%%%%%%%%%%%%%%%%%%%%%%%%%%%%%%%%%%%%%%
	\item[(a)] Write out the lagrangian of the system, and the corresponding equations of motion.
	\resp En el marco del laboratorio,
	\begin{align}
		L &= -mc^2 \sqrt{1 - \frac{\dot{\rv}^2}{c^2}} - \left(\frac{\alpha}{r} + \frac{\beta}{r^2}\right) \nonumber \\
		&= -mc \sqrt{c - \dot{r}^2 - r^2 \dot{\theta}^2 - r^2\dot{\phi}^2 \sin^2 \theta} - \frac{\alpha}{r} - \frac{\beta}{r^2}
	\end{align}
	Usando \texttt{Mathematica}, los momentos generalizados:
	\begin{equation*}
		p_r = \pd{L}{\dot{r}} = \gamma m \dot{r}, ~~~~ p_\theta = \pd{L}{\dot{\theta}} = \gamma m r^2 \dot{\theta}, ~~~~ p_\phi = \pd{L}{\dot{\phi}} = \gamma m r^2 \dot{\phi} \sin^2\theta
	\end{equation*}
	Tenemos las ecuaciones de Euler-Lagrange:
	\begin{equation*}
		\td{p_i}{t} = \pd{L}{q_i}
	\end{equation*}
	\begin{align*}
		\td{}{t} \left[\gamma m \dot{r}\right] &= \pd{}{r} \left[-mc \sqrt{c - \dot{r}^2 - r^2 \dot{\theta}^2 - r^2\dot{\phi}^2 \sin^2 \theta} - \frac{\alpha}{r} - \frac{\beta}{r^2}\right]
	\end{align*}
	
%%%%%%%%%%%%%%%%%%%%%%%%%%%%%%%%%%%%%%%%%%%%%%%%%%%%%%%%%%%%%%%%%%%%%%%%%%%%%%%%%%%%%%%%%%%%%%%%%%%%%%%%
	\item[(b)] Identify all symmetries of the action and construct all integrals of motion which might exist in this problem.
	\resp
	
%%%%%%%%%%%%%%%%%%%%%%%%%%%%%%%%%%%%%%%%%%%%%%%%%%%%%%%%%%%%%%%%%%%%%%%%%%%%%%%%%%%%%%%%%%%%%%%%%%%%%%%%
	\item[(c)] Try to integrate the equations of motion and find trajectories of bound orbits.
	\resp
	
%%%%%%%%%%%%%%%%%%%%%%%%%%%%%%%%%%%%%%%%%%%%%%%%%%%%%%%%%%%%%%%%%%%%%%%%%%%%%%%%%%%%%%%%%%%%%%%%%%%%%%%%
	\item[(d)] Analyze the behavior of your solution in the limit $c\rightarrow\infty$. Assuming that $c$ is large yet finite, find explicitly the first relativistic correction to the trajectory.
	\resp
	
\end{enumerate}

\section{Problema 3}
In the lectures 12-13 we discussed that the covariant equations of motion for the massive particle in the electromagnetic field $A^\mu$ may be obtained from the action
\begin{equation}
	S = \int d\tau \, \Lambda, ~~~~~~~ \Lambda = -mc^2 \sqrt{\dot{x}_\mu \dot{x}^\mu} + e \dot{x}_\mu A^\mu (x) \label{eq 3.1}
\end{equation}
where $\tau$ is the proper time, $x^\mu (\tau)$ is the trajectory of the particle, and $\dot{x}^\mu = \ltd{x^\mu}{\tau}$ is the 4-velocity of the particle. Assume that we replaced the covariant lagrangian \ref{eq 3.1} with
\begin{equation}
	\tilde{\Lambda} = \frac{1}{2} mc^2 \dot{x}_\mu \dot{x}^\mu + e\dot{x}_\mu A^\mu (x) \label{eq 3.2}
\end{equation}

\begin{enumerate}	
%%%%%%%%%%%%%%%%%%%%%%%%%%%%%%%%%%%%%%%%%%%%%%%%%%%%%%%%%%%%%%%%%%%%%%%%%%%%%%%%%%%%%%%%%%%%%%%%%%%%%%%%
	\item[(a)] Demonstrate that the action $S$ is covariant under transformations of the Lorentz group when we use the lagrangian \ref{eq 3.2}.
	\resp $S$ es covariante porque sus componentes ($\tilde{\Lambda}$ y $\tau$) están escritas usando objetos covariantes.
	
	\qed
	
%%%%%%%%%%%%%%%%%%%%%%%%%%%%%%%%%%%%%%%%%%%%%%%%%%%%%%%%%%%%%%%%%%%%%%%%%%%%%%%%%%%%%%%%%%%%%%%%%%%%%%%%
	\item[(b)] Write out the generalized momentum $\mathcal{P}_\mu = \ltd{\tilde{\Lambda}}{\dot{x}^\mu}$ and the equations of motion which correspond to the lagrangian \ref{eq 3.2}.
	\resp El momento
	\begin{align*}
		\mathcal{P}_\mu &= \td{\tilde{\Lambda}}{\dot{x}^\mu}\\
		&= \td{}{\dot{x}^\mu} \left[\frac{1}{2} mc^2 \dot{x}_\alpha \dot{x}^\alpha + e\dot{x}_\gamma A^\gamma (x)\right]\\
		&= \frac{1}{2} mc^2 \td{}{\dot{x}^\mu} \left[\eta_{\alpha\beta} \dot{x}^\beta \dot{x}^\alpha \right] + e A^\gamma (x) \td{}{\dot{x}^\mu} \left[\eta_{\lambda\gamma} \dot{x}^\lambda\right]\\
		&= \frac{1}{2} mc^2 \eta_{\alpha\beta} \left[\delta_\mu^\beta \dot{x}^\alpha + \dot{x}^\beta \delta_\mu^\alpha \right] + e A^\gamma (x) \eta_{\lambda\gamma} \delta_\mu^\lambda\\
		&= \frac{1}{2} mc^2 \left[\eta_{\alpha\mu} \dot{x}^\alpha + \eta_{\mu\beta} \dot{x}^\beta \right] + e A^\gamma (x) \eta_{\mu\gamma}\\
		&= \frac{1}{2} mc^2 \left[\dot{x}_\mu + \dot{x}_\mu \right] + e A_\mu (x)\\
		&= mc^2 \dot{x}_\mu + e A_\mu (x)
	\end{align*}
	Las ecuaciones de Euler-Lagrange:
	\begin{align*}
		\pd{\tilde{\Lambda}}{x^\mu} &= \td{\mathcal{P}_\mu}{\theta}\\
		e\dot{x}_\alpha \partial_\mu A^\alpha (x) &= mc^2 \ddot{x}_\mu + e \td{A_\mu (x)}{\tau}
	\end{align*}
	\begin{align*}
		\Rightarrow mc^2 \ddot{x}_\mu &= e \left(\dot{x}_\alpha \partial_\mu A^\alpha (x) - \pd{A_\mu (x)}{x^\beta} \td{x^\beta}{\tau}\right)\\
		mc^2 \ddot{x}_\mu &= e \left(\eta_{\alpha\lambda} \dot{x}^\lambda \partial_\mu A^\alpha - \dot{x}^\beta \partial_\beta A_\mu \right)\\
		mc^2 \ddot{x}_\mu &= e \left(\dot{x}^\lambda \partial_\mu A_\lambda - \dot{x}^\beta \partial_\beta A_\mu \right)\\
		mc^2 \ddot{x}_\mu &= e \dot{x}^\nu \left(\partial_\mu A_\nu - \partial_\nu A_\mu \right)
	\end{align*}
	\qed
	
%%%%%%%%%%%%%%%%%%%%%%%%%%%%%%%%%%%%%%%%%%%%%%%%%%%%%%%%%%%%%%%%%%%%%%%%%%%%%%%%%%%%%%%%%%%%%%%%%%%%%%%%
	\item[(c)] Demonstrate that your equations of motion are equivalent to the equations of motion found in lectures 12,13 using the lagrangian \ref{eq 3.1}. Demonstrate that the same equations can be obtained with any lagrangian $\bar{\Lambda}$ given by
	\begin{equation}
		\bar{\Lambda} = m c^2 f(\dot{x}_\mu \dot{x}^\mu) + e \dot{x}_\mu A^\mu (x)  \label{eq 3.3}
	\end{equation}
	where $f (y)$ is an arbitrary function which satisfies $\lpd{f(y)}{y} = 1/2$ at $y = 1$ (clearly, the lagrangians (\ref{eq 3.1}, \ref{eq 3.2}) correspond to $f (y) = \sqrt{y}$ and $f (y) = y/2$, respectively).
	\resp Los momentos generalizados del primer lagrangiano serán
	\begin{align*}
		\mathcal{P}_\mu &= \td{\Lambda}{\dot{x}^\mu}\\
		&= \td{}{\dot{x}^\mu} \left[-mc^2 \sqrt{\dot{x}_\mu \dot{x}^\mu} + e \dot{x}_\mu A^\mu (x) \right]\\
		&= -mc^2 \td{}{\dot{x}^\mu} \sqrt{\eta_{\alpha\beta} \dot{x}^\alpha \dot{x}^\beta} + e A^\lambda (x) \td{}{\dot{x}^\mu} \left(\eta_{\lambda\gamma} \dot{x}^\gamma\right)\\
		&= -mc^2 \frac{1}{2 \sqrt{\eta_{\alpha\beta} \dot{x}^\alpha \dot{x}^\beta}} \cdot 2\dot{x}_\mu + e A^\lambda (x) \eta_{\lambda\gamma} \delta_\mu^\gamma\\
		&= -mc^2 \frac{\dot{x}_\mu}{\sqrt{\dot{x}_\alpha \dot{x}^\alpha}} + e A_\mu (x)
	\end{align*}
	Por lo que las ecuaciones de Euler-Lagrange,
	\begin{align*}
		\pd{\Lambda}{x^\mu} &= \td{\mathcal{P}_\mu}{\theta}\\
		e\dot{x}_\alpha \partial_\mu A^\alpha (x) &= mc^2 \left(\frac{(\dot{x}_\alpha \dot{x}^\alpha) \ddot{x}_\mu - \dot{x}_\mu (\ddot{x}_\beta \dot{x}^\beta + \dot{x}_\gamma \ddot{x}^\gamma) }{\dot{x}_\alpha \dot{x}^\alpha}\right) + e \td{A_\mu (x)}{\tau}\\
		e\dot{x}_\alpha \partial_\mu A^\alpha (x) &= mc^2 \left(\frac{(\dot{x}_\alpha \dot{x}^\alpha) \ddot{x}_\mu - 2 \dot{x}_\mu (\ddot{x}_\beta \dot{x}^\beta) }{\dot{x}_\alpha \dot{x}^\alpha}\right) + e \td{A_\mu (x)}{\tau}
	\end{align*}
	Pero $\ddot{x}_\beta \dot{x}^\beta = 0$ y $\dot{x}_\alpha \dot{x}^\alpha = c^2$
	\begin{align*}
		e\dot{x}_\alpha \partial_\mu A^\alpha (x) &= mc^2 \left(\frac{c^2 \ddot{x}_\mu - 0}{c^2}\right) + e \td{A_\mu (x)}{\tau}\\
		e\dot{x}_\alpha \partial_\mu A^\alpha (x) &= mc^2 \ddot{x}_\mu + e \td{A_\mu (x)}{\tau}
	\end{align*}
	Lo cual es el mismo resultado encontrado para el caso anterior, por lo que se sigue que las ecuaciones de movimiento son equivalentes.
	
	\qed
	
	Para el caso más general, los momentos generalizados
		\begin{align*}
		\mathcal{P}_\mu &= \td{\bar{\Lambda}}{\dot{x}^\mu}\\
		&= \td{}{\dot{x}^\mu} \left[m c^2 f(\dot{x}_\alpha \dot{x}^\alpha) + e \dot{x}_\mu A^\mu (x) \right]\\
		&= m c^2 \td{}{\dot{x}^\mu} f(\dot{x}_\alpha \dot{x}^\alpha) + e A_\mu (x)\\
		&= m c^2 \td{\left(\dot{x}_\alpha \dot{x}^\alpha\right)}{\dot{x}^\mu} \td{}{\left(\dot{x}_\alpha \dot{x}^\alpha\right)} f(\dot{x}_\alpha \dot{x}^\alpha) + e A_\mu (x)\\
		&= m c^2 \cdot 2\dot{x}_\mu \td{}{\left(\dot{x}_\alpha \dot{x}^\alpha\right)} f(\dot{x}_\alpha \dot{x}^\alpha) + e A_\mu (x)
	\end{align*}
	Ya que $\lpd{f(y)}{y} = 1/2$ en $y = 1$, y $\dot{x}_\alpha \dot{x}^\alpha = c^2 = 1$ (en unidades naturales), tendremos
	\begin{align*}
		\mathcal{P}_\mu &= m c^2 \cdot 2\dot{x}_\mu \cdot \frac{1}{2} + e A_\mu (x)\\
		&= m c^2 \dot{x}_\mu + e A_\mu (x)
	\end{align*}
	lo cual es el mismo resultado para el caso de $\tilde{\Lambda}$ (y de $\Lambda$ si se hacía el reemplazo antes). Ya que el momento no depende de la posición $x_\mu$ y los momentos son idénticos, tendremos que los resultados de las ecuaciones de Euler-Lagrange serán los mismos.
	
	\qed
	
%%%%%%%%%%%%%%%%%%%%%%%%%%%%%%%%%%%%%%%%%%%%%%%%%%%%%%%%%%%%%%%%%%%%%%%%%%%%%%%%%%%%%%%%%%%%%%%%%%%%%%%%
	\item[(d)] Find the nonrelativistic lagrangian $L$ which corresponds to the lagrangian \ref{eq 3.2}. Pay attention that the nonrelativistic lagrangian $L$ contributes to the action as $S = \int L \, dt$.
	\resp Tomando unidades naturales $c = 1$, y considerando que la 4-velocidad se relaciona con la velocidad como $\dot{x}_\mu = \gamma v_\mu$, la acción será
	\begin{align*}
		S &= \int L \, dt\\
		&= \int \tilde{\Lambda} \, d\tau \\
		&= \int \tilde{\Lambda} \frac{1}{\gamma} \, dt
	\end{align*}
	Por lo que identificamos el lagrangiano:
	\begin{align*}
		\Rightarrow L &= \frac{\tilde{\Lambda}}{\gamma}\\
		&= \frac{1}{\gamma} \left(\frac{1}{2} m \dot{x}_\mu \dot{x}^\mu + e\dot{x}_\mu A^\mu (x) \right)\\
		&= \frac{1}{\gamma} \left(\frac{1}{2} m \gamma^2 v_\mu v^\mu + e \gamma v_\mu A^\mu (x) \right)\\
		&= \frac{1}{2} m \gamma v_\mu v^\mu + e v_\mu A^\mu (x)
	\end{align*}
	Ya que $x^\mu = (t, \n{x})$ (donde $\n{x}$ es la posición espacial), entonces $v^\mu = (1, \n{v})$ y $v_\mu = (1, -\n{v})$ (donde $\n{v}$ es la velocidad espacial), tendremos que
	\begin{equation*}
		v_\mu v^\mu = 1 - |\n{v}|^2 = \frac{1}{\gamma^2}
	\end{equation*}
	entonces, considerando $A^\mu = (\varphi, \n{A})$
	\begin{align*}
		L &= \frac{1}{2} m \frac{1}{\gamma} + e v_\mu A^\mu (x)\\
		&= \frac{1}{2} m \sqrt{1 - |\n{v}|^2} + e \left(\varphi (x) + \n{v} \cdot \n{A} (x)\right)
	\end{align*}
	Expandiendo para $\n{v} \rightarrow 0$ utilizando \texttt{Mathematica},
	\begin{equation*}
		L_\mathrm{NR} \approx \frac{1}{2} m - \frac{1}{4} m |\n{v}|^2 + e \left(\varphi (x) + \n{v} \cdot \n{A} (x)\right)
	\end{equation*}
	\begin{equation*}
		????????
	\end{equation*}
	\textbf{Error:} Al parecer no se puede hacer esto con $\tilde{\Lambda}$ por algo que se llama \textit{reparameterization invariance}, si lo hubiera hecho con el $\Lambda$ original me hubiera dado bien. warever\ldots
\end{enumerate}

\section{Problema 4}
A pointlike nonrelativistic particle of mass $m$ scatters in the time-dependent spherically symmetric Coulomb potential
\begin{equation}
	U(r,t) = \frac{U_0(t)}{r}, ~~~~~~~ U_0(t) = \frac{\kappa}{\sqrt{t}}, ~~~~~~~ \kappa > 0 \label{eq 4.1}
\end{equation}
\begin{equation}
	\Rightarrow U(r,t) = \frac{\kappa}{r \sqrt{t}}
\end{equation}

\begin{enumerate}
%%%%%%%%%%%%%%%%%%%%%%%%%%%%%%%%%%%%%%%%%%%%%%%%%%%%%%%%%%%%%%%%%%%%%%%%%%%%%%%%%%%%%%%%%%%%%%%%%%%%%%%%
	\item[(a)] Identify all continuous symmetries of the action. Analyze if there are symmetries of type
	\begin{equation}
		r \rightarrow \alpha r, ~~~ t \rightarrow \beta t, ~~~~~~~ \alpha, \beta = \mathrm{const.}
	\end{equation}
	which do not change the action or the equations of motion.
	\resp Tenemos el lagrangiano
	\begin{equation}
		L = \frac{1}{2}m \left(\dot{r}^2 + r^2 \dot{\theta}^2 + r^2\dot{\phi}^2 \sin^2 \theta \right) - \frac{\kappa}{r \sqrt{t}}
	\end{equation}
	La variable $\phi$ es cíclica, por lo que $M_z$ será conservado. Aun más, el sistema es esféricamente simétrico, por lo que todo el vector de momento angular se conservará:
	\begin{equation}
		\n{M} = \mathrm{const.}
	\end{equation}
	
	Reemplazando $\rv \rightarrow \alpha \rv, t\rightarrow \beta t$ en las derivadas,
	\begin{equation}
		\pd{r}{t} \rightarrow \pd{\alpha r}{\beta t} = \frac{\alpha}{\beta} \pd{r}{t}
	\end{equation}
	Luego en el lagrangiano
	\begin{align*}
		L' &= \frac{1}{2}m \left(\frac{\alpha^2}{\beta^2} \dot{r}^2 + \frac{\alpha^2}{\beta^2} r^2 \dot{\theta}^2 + \frac{\alpha^2}{\beta^2} r^2\dot{\phi}^2 \sin^2 \theta \right) - \frac{\kappa}{\alpha r \cdot \sqrt{\beta} \sqrt{t}} \\
		L' &= \frac{\alpha^2}{\beta^2} \left[\frac{1}{2}m \left(\dot{r}^2 + r^2 \dot{\theta}^2 + r^2\dot{\phi}^2 \sin^2 \theta \right)\right] - \frac{1}{\alpha \sqrt{\beta}} \frac{\kappa}{r \sqrt{t}}
	\end{align*}
	Ya que multiplicar un lagrangiano por una cantidad constante no cambia las ecuaciones de movimiento, tendremos una simetría sí y solo sí $\alpha = \sqrt{\beta}$.
	
	Sin embargo, al considerar el efecto sobre la acción,
	\begin{align*}
		S = \int L \, dt ~\rightarrow~ S' &= \int L' \, dt' \\
		&= \int \left(\frac{\alpha^2}{\beta^2} \left[\frac{1}{2}m \left(\dot{r}^2 + r^2 \dot{\theta}^2 + r^2\dot{\phi}^2 \sin^2 \theta \right)\right] - \frac{1}{\alpha \sqrt{\beta}} \frac{\kappa}{r \sqrt{t}}\right) \, \beta dt\\
		&= \int \left(\frac{\alpha^2}{\beta} \left[\frac{1}{2}m \left(\dot{r}^2 + r^2 \dot{\theta}^2 + r^2\dot{\phi}^2 \sin^2 \theta \right)\right] - \frac{\sqrt{\beta}}{\alpha} \frac{\kappa}{r \sqrt{t}}\right) \, dt\\
	\end{align*}
	Para que podamos usar el teorema de Noether la acción deberá ser invariante, por lo que deberá cumplirse simultáneamente
	\begin{equation*}
		\frac{\alpha^2}{\beta} = 1, ~~~~~~~ \frac{\sqrt{\beta}}{\alpha} = 1
	\end{equation*}
	\begin{equation*}
		\Rightarrow \beta = \alpha^2
	\end{equation*}
	\begin{equation*}
		\Rightarrow \rv \rightarrow \alpha \rv, ~~~~~ t\rightarrow \alpha^2 t, ~~~~~ \alpha = \mathrm{const.}
	\end{equation*}
	\qed
	
%%%%%%%%%%%%%%%%%%%%%%%%%%%%%%%%%%%%%%%%%%%%%%%%%%%%%%%%%%%%%%%%%%%%%%%%%%%%%%%%%%%%%%%%%%%%%%%%%%%%%%%%
	\item[(b)] Construct the integrals of motion which might exist according to Noether’s theorem. Demonstrate that these quantities are conserved using the equations of motion.
	\resp El momento angular vendrá dado por
	\begin{equation}
		\n{M} = m\rv\times\dot{\rv} = \mathrm{const.}
	\end{equation}
	
	Para la simetría de escala, escribiendo las transformaciones en notación con $\varepsilon$
	\begin{equation*}
		\alpha = e^{\varepsilon(t)} \approx 1 + \varepsilon(t)
	\end{equation*}
	tal que $\alpha$ es generado por cambios infinitesimales $\varepsilon$, los cambios de coordenadas y tiempo será, entonces
	\begin{equation*}
		\rv' \approx \rv + \delta \rv \approx \rv + \varepsilon \, \rv, ~~~~~ t' \approx t + \delta t \approx t + (1 + \varepsilon)^2 t \approx t + 2\varepsilon \, t
	\end{equation*}
	\begin{equation*}
		\delta \rv \approx \varepsilon \rv, ~~~~~ \delta t \approx 2\varepsilon t
	\end{equation*}
	\begin{align*}
		\Rightarrow d \rv' &= d \left(e^{\varepsilon(t)} \rv\right)\\
		&= e^{\varepsilon} d\rv + \rv e^{\varepsilon} d \varepsilon\\
		&= e^{\varepsilon} \dot{\rv} dt + \rv e^{\varepsilon} \dot{\varepsilon} dt \\
		&\approx (1 + \varepsilon) \dot{\rv} dt + \rv (1 + \varepsilon) \dot{\varepsilon} dt \\
		&\approx \underbrace{\dot{\rv} dt}_{d\rv} + \underbrace{\varepsilon \dot{\rv} dt}_{\mathcal{O}(\varepsilon)} + \underbrace{\rv \dot{\varepsilon} dt}_{\mathcal{O}(\dot{\varepsilon})} + \underbrace{\rv \varepsilon \dot{\varepsilon} dt}_{\mathcal{O}(\varepsilon \dot{\varepsilon})}\\
		&\\
		\Rightarrow d t' &= d \left(e^{2\varepsilon(t)} t\right)\\
		&= e^{2\varepsilon} dt + t e^{2\varepsilon} \cdot 2 d \varepsilon\\
		&= e^{2\varepsilon} dt + 2 t e^{2\varepsilon} \dot{\varepsilon} dt \\
		&\approx (1 + \varepsilon)^2 dt + 2 t (1 + \varepsilon)^2 \dot{\varepsilon} dt \\
		&\approx (1 + 2\varepsilon + \varepsilon^2) dt + 2 t (1 + 2\varepsilon + \varepsilon^2) \dot{\varepsilon} dt \\
		&\approx dt + \underbrace{2\varepsilon dt}_{\mathcal{O}(\varepsilon)} + \underbrace{\varepsilon^2 dt}_{\mathcal{O}(\varepsilon^2)} + \underbrace{2 t \dot{\varepsilon} dt}_{\mathcal{O}(\dot{\varepsilon})} + \underbrace{4 t \varepsilon \dot{\varepsilon} dt}_{\mathcal{O}(\varepsilon \dot{\varepsilon})} + \underbrace{2 t \varepsilon^2 \dot{\varepsilon} dt}_{\mathcal{O}(\varepsilon^2 \dot{\varepsilon})}
	\end{align*}
	Todos los términos $\varepsilon \dot{\varepsilon}, ~ \varepsilon^2, ~ \varepsilon^2 \dot{\varepsilon}$ son insignificantes, y después de esto deberemos diferenciar por $\dot{\varepsilon}$, por lo que solo nos interesarán los términos $\mathcal{O}(\dot{\varepsilon})$
	\begin{equation*}
		d \rv' \approx \left(\dot{\rv} + \dot{\varepsilon} \rv\right) dt, ~~~~~ dt' \approx \left(1 + 2 \dot{\varepsilon} t\right) dt
	\end{equation*}
	Luego, la velocidad
	\begin{equation*}
		\dot{\rv}' \equiv \frac{d \rv'}{d t'} \approx \frac{\left(\dot{\rv} + \dot{\varepsilon} \rv \right) dt}{\left(1 + 2 \dot{\varepsilon} t\right) dt}
	\end{equation*}
	Aproximando la división con \texttt{Mathematica},
	\begin{equation*}
		\dot{\rv}' \approx \dot{\rv} + \delta \dot{\rv} \approx \dot{\rv} + \dot{\varepsilon}\left(\rv - 2 \dot{\rv} t\right)
	\end{equation*}
	\begin{equation*}
		\delta \dot{\rv} \approx \dot{\varepsilon}\left(\rv - 2 \dot{\rv} t\right)
	\end{equation*}
	Luego, aproximando el lagrangiano a primer orden,
	\begin{align*}
		L'(\rv',\dot{\rv}',t') &= L'(\rv + \delta \rv, \dot{\rv} + \delta \dot{\rv}, t + \delta t)\\
		&\approx L(\rv,\dot{\rv},t) + \pd{L}{\rv} \cdot \delta\rv + \pd{L}{\dot{\rv}} \cdot \delta\dot{\rv} + \pd{L}{t} \delta t\\
		&\approx L(\rv,\dot{\rv},t) + \pd{L}{\rv} \cdot \varepsilon \rv + \pd{L}{\dot{\rv}} \cdot \dot{\varepsilon}\left(\rv - 2 \dot{\rv} t\right) + \pd{L}{t} 2 \varepsilon t
	\end{align*}
	De nuevo, solo nos interesarán los términos $\mathcal{O}(\dot{\varepsilon})$, por lo que, considerando que $\lpd{L}{\dot{\rv}} = m \dot{\rv}$
	\begin{equation*}
		L' \approx L + m \dot{\rv} \cdot \dot{\varepsilon}\left(\rv - 2 \dot{\rv} t\right) 
	\end{equation*}
	Por lo que la acción
	\begin{align*}
		S' &= \int L' \, dt'\\
		&\approx \int \left[L + m \dot{\rv} \cdot \dot{\varepsilon}\left(\rv - 2 \dot{\rv} t\right)\right] \, \left[1 + 2 \dot{\varepsilon} t\right] dt\\
		&\approx \int L \, dt + \int m \dot{\rv} \cdot \dot{\varepsilon} \left(\rv - 2 \dot{\rv} t\right) \, dt + \int 2 \dot{\varepsilon} t L \, dt + \int 2 m \dot{\rv} \cdot \dot{\varepsilon}^2 t \left(\rv - 2 \dot{\rv} t\right) \, dt
	\end{align*}
	El último término es $\mathcal{O}(\dot{\varepsilon}^2)$, por lo que es insignificante. Tendremos entonces
	\begin{equation*}
		S' \approx \int \left[L + m \dot{\rv} \cdot \dot{\varepsilon} \left(\rv - 2 \dot{\rv} t\right) + 2 \dot{\varepsilon} t L\right] \, dt
	\end{equation*}
	Y la variación de $L$ será
	\begin{equation*}
		S' \approx \int \left[L + \delta L\right] dt \approx \int \left[L + \dot{\varepsilon} \left( m \dot{\rv} \cdot \left(\rv - 2 \dot{\rv} t\right) + 2 t L\right)\right] \, dt
	\end{equation*}
	\begin{equation*}
		\delta L \approx \dot{\varepsilon}\left[ m \dot{\rv} \cdot \left(\rv - 2 \dot{\rv} t\right) + 2 t L\right]
	\end{equation*}
	Finalmente, la cantidad conservada será
	\begin{align*}
		Q &= \pd{\delta L}{\dot{\varepsilon}}\\
		&= m \dot{\rv} \cdot \left(\rv - 2 \dot{\rv} t\right) + 2 t L\\
		&= \pv \cdot \rv - 2 m \dot{\rv}^2 t + 2 t \left(\frac{1}{2}m \dot{\rv}^2 - U\right)\\
		&= \pv \cdot \rv - 2 t \left(\frac{1}{2}m \dot{\rv}^2 + U\right)\\
		&= \pv \cdot \rv - 2 H t\\
	\end{align*}
	
%%%%%%%%%%%%%%%%%%%%%%%%%%%%%%%%%%%%%%%%%%%%%%%%%%%%%%%%%%%%%%%%%%%%%%%%%%%%%%%%%%%%%%%%%%%%%%%%%%%%%%%%
	\item[(c)] We have seen in lectures that in static (time-independent) Coulomb potential the particle's trajectory always remains in the plane (so-called scattering plane). Analyze if this property remains valid in the time-dependent potential given in Eq. \ref{eq 4.1}. Analyze if the Kepler’s laws remain valid in this potential field.
	\resp El potencial es central y el momento angular $\n{M}$ es conservado componente por componente, por lo que el plano de rotación será constante.
	
	La segunda ley de Kepler se conservará porque el campo es central; la primera y tercera no lo harán, ya que dependen de una trayectoria específica que en este caso será alterada.
	
	\qed
\end{enumerate}

\section{Problema 5}
A relativistic particle of mass $m$ moves in the spherically symmetric potential $U (r) = kr^2/2$ of 3D harmonic oscillator.
\begin{enumerate}
%%%%%%%%%%%%%%%%%%%%%%%%%%%%%%%%%%%%%%%%%%%%%%%%%%%%%%%%%%%%%%%%%%%%%%%%%%%%%%%%%%%%%%%%%%%%%%%%%%%%%%%%
	\item[(a)] Write out the lagrangian of the system, and the corresponding equations of motion
	\resp El lagrangiano relativista en el marco del potencial
	\begin{equation}
		L = -m c \sqrt{c^2 - \left(\dot{r}^2 + r^2 \dot{\theta}^2 + r^2\dot{\phi}^2 \sin^2 \theta \right)} - \frac{1}{2} k r^2
	\end{equation}
	los momentos generalizados, usando \texttt{Mathematica},
	\begin{equation}
		p_r = \gamma m \dot{r}, ~~~~~ p_\theta = \gamma m r^2 \dot{\theta}, ~~~~~ p_\phi = \gamma m r^2 \dot{\phi} \sin^2\theta
	\end{equation}
	y las ecuaciones de Euler-Lagrange:
	\begin{align*}
		\td{}{t} \left( \gamma m \dot{r} \right) &= \frac{1}{2} \gamma m \left(2 r \dot{\phi}^2 \sin ^2 \theta + 2 r \dot{\theta}^2\right) - k r\\
		\td{}{t} \left( \gamma m r^2 \dot{\theta} \right) &= \frac{1}{2} \gamma m r^2 \dot{\phi}^2 \sin 2\theta\\
		\td{}{t} \left( \gamma m r^2 \dot{\phi} \sin^2\theta \right) &= 0
	\end{align*}
	De aquí, $p_\phi = \mathrm{const.}$, por lo que
	\begin{align*}
		\td{}{t} \left( \gamma m \dot{r} \right) &= \frac{p_\phi^2}{\gamma m r^2 \sin^2\theta} + \gamma m r \dot{\theta}^2 - k r\\
		\td{}{t} \left( \gamma m r^2 \dot{\theta} \right) &= \frac{1}{2} m \gamma r^2 \dot{\phi}^2 \sin 2\theta\\
		\td{}{t} \left( \gamma m r^2 \dot{\phi} \sin^2\theta \right) &= 0
	\end{align*}
	
	\qed
	
%%%%%%%%%%%%%%%%%%%%%%%%%%%%%%%%%%%%%%%%%%%%%%%%%%%%%%%%%%%%%%%%%%%%%%%%%%%%%%%%%%%%%%%%%%%%%%%%%%%%%%%%
	\item[(b)] Identify all the symmetries of the action and construct all integrals of motion which might exist in this problem.
	\resp
	
	\qed
	
%%%%%%%%%%%%%%%%%%%%%%%%%%%%%%%%%%%%%%%%%%%%%%%%%%%%%%%%%%%%%%%%%%%%%%%%%%%%%%%%%%%%%%%%%%%%%%%%%%%%%%%%
	\item[(c)] Try to integrate the equations of motion and find trajectories for bound orbits. Analyze if the trajectories are closed for bound orbits and if the motion remains periodic as a function of time.
	\resp
	
	\qed
	
%%%%%%%%%%%%%%%%%%%%%%%%%%%%%%%%%%%%%%%%%%%%%%%%%%%%%%%%%%%%%%%%%%%%%%%%%%%%%%%%%%%%%%%%%%%%%%%%%%%%%%%%
	\item[(d)] Demonstrate that in the limit $c\rightarrow\infty$ your results reproduce well-known nonrelativistic limit. Assuming that $c$ is large yet finite, find explicitly the first relativistic correction to trajectory.
	\resp
	
	\qed
	
\end{enumerate}








\end{document}